%% ma: L12-B07-0007
%% lop: 12
%% chuong: 2
%% bai: 7
%% dang: 12.7.6
%% loai: TN4
%% muc: NB
%% dap_an: "A"
%% nguon: "LaTeX giáo viên gửi 21/10/2026 (M05, Vectơ và hệ trục tọa độ trong không gian, 107 câu), Phần I Câu 23"
%% kien_thuc: "Bài 7 (tọa độ điểm, hình chiếu lên mặt phẳng tọa độ), Bài 8 (trung điểm)"
%% trang_thai: chinh-thuc
%% ghi_chu: "Giáo viên đã duyệt 21/10/2026."
\begin{ex}
Trong không gian $Oxyz$, điểm đối xứng với điểm $A(1;-3;5)$ qua mặt phẳng $(Oyz)$ có tọa độ là
\choice{\True $(-1;-3;5)$}{$(1;3;-5)$}{$(-1;3;-5)$}{$(1;-3;-5)$}
\loigiai{Hình chiếu của $A$ lên $(Oyz)$ là $H(0;-3;5)$, $H$ là trung điểm của $AA'$ nên $A'(-1;-3;5)$.}
\end{ex}
